\documentclass[12pt]{article}
\usepackage[T1]{fontenc}
\usepackage[utf8]{inputenc}
\usepackage{lmodern}
\usepackage{textcomp}
\usepackage{enumitem}
\usepackage{geometry}
\usepackage{graphicx}
\usepackage{amsmath, amssymb}
\geometry{margin=1in}
\begin{document}
\begin{center}
Constitutional Law - Undergraduate Level Assessment 20 \
Total Marks: 40 \
Answer all questions.
\end{center}

\section*{Multiple Choice (10 marks)}
\begin{enumerate}[label=\arabic*.]
\item Dworkin contends that to every legal question there is only one right answer. Which proposition below is most inconsistent with this claim?
\begin{enumerate}[label=(\alph*)]
    \item In hard cases judges generally decide cases on the basis of rights.
    \item The rights of the parties feature in the determination of most cases before the courts.
    \item Judges exercise strong discretion.
    \item Judges seek the best 'fit' with constitutional and institutional history.
\end{enumerate}
\item Which one of the following doctrines is not associated with Natural Law thinking
\begin{enumerate}[label=(\alph*)]
    \item Doctrine of substituted security
    \item Doctrine of due process
    \item Doctrine of bias
    \item Doctrine of reasonableness
\end{enumerate}
\item Which statement below best represents Durkheim's view of the function of punishment?
\begin{enumerate}[label=(\alph*)]
    \item Deterrence.
    \item Rehabilitation.
    \item Vengeance.
    \item Desert.
\end{enumerate}
\item Which statement below is the least likely to follow logically from Savigny's notion of a Volksgeist?
\begin{enumerate}[label=(\alph*)]
    \item A society's law is a reflection of its culture.
    \item Law is like language.
    \item Law is the deliberate expression of a sovereign's will.
    \item Law is an integral element of the social fabric.
\end{enumerate}
\item Robert Nozick's proposes the formula r x H as a guide to determine the appropriate punishment. What does it mean?
\begin{enumerate}[label=(\alph*)]
    \item Effectiveness of rehabilitation of multiplied by hazard to the community.
    \item Extent of responsibility multiplied by actual harm done.
    \item Risk of violence multiplied by degree of humility of offender.
    \item Recidivism multiplied by defendant's history.
\end{enumerate}
\end{enumerate}

\section*{True / False (10 marks)}
\textit{Answer True or False}
\begin{enumerate}[label=\arabic*.]
\setcounter{enumi}{5}
\item True or False: Kelsen's Grundnorm is a presupposition that facilitates our understanding of a legal system.
\item True or False: Standard-setting in human rights diplomacy means establishing binding legal standards that countries must follow.
\item True or False: Moral values are independent and objective, playing a crucial role in understanding law and justice.
\item True or False: The UN Human Rights Council is a treaty-based mechanism that evaluates the human rights records of all UN member states.
\item True or False: The flag State ordinarily exercises jurisdiction over crimes committed on board its vessels.
\end{enumerate}

\section*{Long Form Response (20 marks)}
\textit{Answer each question with a clear and complete explanation.}
\begin{enumerate}[label=\arabic*.]
\setcounter{enumi}{10}
\item Discuss the role of treaties as sources of international law under Article 38 of the ICJ Statute, emphasizing the significance of treaties that are in force and binding on the parties involved.
\item Discuss the rationale behind the exhaustion of local remedies in international human rights law, including its implications for local judicial capacity and international tribunal efficiency.
\end{enumerate}

\vfill
\noindent\textit{End of Paper}
\end{document}